\documentclass[output=paper, 
colorlinks,
citecolor=brown,
newtxmath
]{langscibook} 
\bibliography{localbibliography}

\author{Rafael Ignacio Zurita\affiliation{Universidad Nacional del Comahue}}
\title{Dispositivos de E/S (periféricos)}  
\abstract{
Además del procesador y memoria, la mayoría de los sistemas embebidos
contienen varios dispositivos de hardware extra. Algunos de estos
dispositivos son específicos a la aplicación, mientras que otros -como los
relojes y puertos seriales- son útiles de una manera general a una gran 
variedad de sistemas distintos.
Los que residen junto con el procesador dentro del mismo chip se los
denomina periféricos internos, u on-chip. Los que se encuentran fuera
del chip donde está el procesador se los denomina de manera opuesta,
es decir, periféricos externos. En este capítulo se detalla la
mayoría de los problemas de software comunes que surgen cuando se programan
periféricos de ambos tipos.

\keywords{sistema embebido, dispositivos de E/S, periféricos, registro de E/S, controladores de E/S, drivers de dispositivos}
}


\input{localpackages.tex}
\input{localcommands.tex}


\usepackage[spanish]{babel}
\usepackage[bottom]{footmisc}
\usepackage{tabularx}

\definecolor{aliceblue}{rgb}{0.90, 0.93, 0.96}


\begin{document}

\selectlanguage{spanish}

\chapterfont{\Large\color{LightBlue}} 
\chapter*{Programación de Sistemas Embebidos 2020\\ Dispositivos de E/S (periféricos)}

\begingroup
\let\clearpage\relax
\cleardoublepage
\hypersetup{linkcolor=blue}
\tableofcontents*
\endgroup

{\def\addcontentsline#1#2#3{}\maketitle}

\setcounter{page}{1}


\hfill\begin{minipage}{0.8\linewidth} \footnotesize
The moral is obvious. You can't trust code that you did
not totally create yourself. (Especially code from companies
that employ people like me.) No amount of
source-level verification or scrutiny will protect you
from using untrusted code. In demonstrating the possibility
of this kind of attack, I picked on the C compiler.
I could have picked on any program-handling program
such as an assembler, a loader, or even hardware microcode.
As the level of program gets lower, these bugs
will be harder and harder to detect. A well-installed
microcode bug will be almost impossible to detect.\\ 
—Ken Thompson aug-1984, Reflections on Trusting Trust.
\end{minipage}

% \togglepaper[0]

\def\maketitle{

% Titulo 
 \makeatletter
 {\color{bl} \centering \huge \sc \textbf{
\large \vspace*{-8pt} \color{black} Programación de Sistemas Embebidos
 \vspace*{8pt} }\par}
 \makeatother


% Autor
 \makeatletter
 {\centering \small 
 	Departamento de Ingeniería de Computadoras \\
 	Facultad de Informática - Universidad Nacional del Comahue \\
 	\vspace{20pt} }
 \makeatother

}









\section {Sistemas de Tiempo Real y RTOS}


\subsection{Sistemas de tiempo real}
Un sistema es de tiempo real cuando en la especificación de sus requerimientos
existen requisitos de tiempo de respuesta limites.
Por lo que su correctitud, una vez implementado, esta definida por la correctitud de los resultados computados, y porque cumple con los tiempos de respuesta definidos en sus especificaciones. Si el sistema no es capaz de responder a un evento o tarea —definida en los requerimientos del sistema— en el tiempo maximo permitido, tambien definido en las especificaciones, entonces se dice que el sistema falló, y no es de tiempo real.

Se los categoriza en dos tipo. \textbf{Hard Real-Time (tiempo real estricto)}: Cuando los plazos temporales deben ser cumplidos siempre. En este tipo están los sistemas que ante un fallo corre riesgo la vida humana o la estabilidad del sistema. Por ejemplo, la activación de un airbag en un automóvil.  \textbf{Soft Real-Time (tiempo real flexible)}: Son los sistemas en donde incumplir ocasionalmente un tiempo de respuesta degrada el rendimiento, pero no provoca un fallo crítico. Por ejemplo, un sistema de reproducción de video en streaming.

\begin{center}
\begin{tabularx}{\textwidth}{|X|}
%\hline
\rowcolor{aliceblue}
\textbf{Un Sistema es de Tiempo Real si..}\\
  1.  Los resultados computados son correctos. \\
  2.  Se cumplen los tiempos de respuesta definidos en sus especificaciones. \\
Un resultado producido fuera del plazo definido equivale, en este contexto, a un sistema que NO es de tiempo real.
%\hline
\end{tabularx}
\end{center}


Usualmente, al diseñar un sistema de tiempo real, se definen dos clases de tareas que son llevadas a cabo: tareas periódicas estrictas, y tareas activadas por eventos. Para cada uno de estas tareas se define tambien, con precision, el tiempo de respuesta maximo permitido por parte del sistema (sobre todo para las tareas activadas por eventos). Luego, este diseño se divide usualmente en tareas independientes, que seran ejecutadas de manera concurrente en tiempo de ejecucion. 

Como software de apoyo, se utiliza un sistema operativo de tiempo real 
(RTOS, por su sigla en ingles en Real Time Operating System).
Estos se ejecutan principalmente en microprocesadores de baja complejidad y, mayormente, en microcontroladores. 
Este software es el encargado de ejecutar las tareas de manera concurrente, y tambien da soporte para lograr los tiempos de respuesta requeridos.
Sin embargo, esto no es un proceso automatico. Es decir, utilizar un sistema 
operativo de tiempo real no implica que el sistema logrado sea de tiempo real. 


\begin{center}
\begin{tabularx}{\textwidth}{|X|}
%\hline
\rowcolor{aliceblue}
\textbf{Importante}\\
	Utilizar un RTOS NO garantiza automaticamente que el sistema sea de tiempo real. Para lograrlo, los desarrolladores deben diseñarlo apropiadamente. Por ejemplo, asignando prioridades correctas a cada tarea (si el sistema gestiona tareas con prioridades), asignando eventos a cada tarea, y definiendo como deben sincronizarse y comunicarse entre si.
%\hline
\end{tabularx}
\end{center}


\subsection{Sistema Operativo de Tiempo Real}

Existen distintos metodos y mecanismos que un sosftware de apoyo RTOS puede incluir para el desarrollo de un sistema de tiempo real.
Uno ampliamente utilizado es el uso de un RTOS con planificador de prioridad fija y apropiativo. Cuando el sistema está en ejecucion, cada vez que el RTOS activa una tarea, esto es, la coloca en estado de "listo para ejecutar", verifica su prioridad. Si esta prioridad es más alta que la de la tarea que actualmente utiliza la CPU, entonces el RTOS realiza un cambio de contexto y coloca a la nueva tarea de mas alta prioridad a ejecutar.

Si existen varias tareas listas para ejecutar, con la misma prioridad que la tarea que esta actualmente utilizando la CPU, entonces el RTOS realiza una expropiacion de la CPU en intervalos regulares, empleando la planificación round-robin para asignarla a
otra tarea. Aquí, cada tarea se ejecuta durante un QUANTUM de tiempo fijo, antes de ser interrumpida, para que la CPU sea luego asignada a otra tarea de la misma prioridad. Luego de que todas las tareas de la misma prioridad se ejecutaron
durante su QUANTUM, el RTOS volverá a asignar la CPU a la primera tarea
de esta cola FIFO.


Contando con estas caracteristicas, los desarrolladores pueden diseñar un sistema cuyo funcionamiento es predecible, lo cual es una condicion necesaria para evaluar y conocer si el sistema resultante cumple con los plazos definidos y puede, por lo tanto, ser de tiempo real.

En el momento de modelado y diseño del sistema, con esta técnica en mente, 
el desarrollador 
debe especificar para cada tarea una prioridad, de acuerdo a si esa tarea
debe tener o no un tiempo de respuesta limite. Tambien relacionar las tareas 
con los eventos, y definir como deben sincronizarse y comunicarse esas tareas. 
En tiempo de ejecucion, el sistema operativo de tiempo real será el encargado 
de asegurar que siempre se este ejecutando la tarea de mayor prioridad. 

Como el comportamiento del sistema bajo un RTOS es predecible
el desarrollador del sistema embebido puede evaluar el modelo diseñado 
ante todas las situaciones posibles, para corroborar que siempre 
se cumplen los tiempos de respuestas definidos.
Si la cantidad de estados del sistema y situaciones posibles son 
inabarcables, los desarrolladores pueden realizar simulaciones y 
aproximaciones para evaluar si el sistema cumplirá con los 
plazos requeridos.

Existe tambien un modelo formal para validar un modelo de sistema de tiempo real
el cual se basa en convertir los eventos asincronicos a periodicos, y evaluar
todo el sistema con por ejemplo ............
Puede leer mayor documentacion en : AGREGAR REFERENCIA.

\subsubsection{Caracteristicas de un RTOS}

En sintesis, un RTOS que implementa el esquema de prioridad fija y 
apropiatividad (que es el caso más común) debe contar con las siguientes 
caracteristicas:
\begin{itemize}
	\item Administrador de tareas: debe poder ejecutar tareas diferentes 
	via multiprogramacion o multi-hilos,  y ser apropiativo (preemptive).
\item Planificador por prioridades: cada tarea debe contar con una prioridad, y el RTOS se encarga de que la tarea de mayor prioridad es la que utiliza la CPU.
\item Sincronizacion y comunicacion entre tareas: el RTOS debe contar con mecanismos para que las tareas puedan coordinarse y transferirse datos de forma segura (semaforos, mutex, colas de mensajes, etc.).
\item Herencia de prioridad: debe existir un mecanismo de herencia de prioridad para evitar la inversion de prioridades, esto es, que una tarea de baja prioridad no bloquee indefinidamente a una de alta prioridad.
\item Eventos asincronicos: el RTOS debe permitir eventos asincronicos, como pueden ser las interrupciones de E/S, permitiendo reaccionar ante eventos externos sin consulta activa (polling).
\end{itemize}


\subsubsection{ Donde se usan los RTOS}

Los sistemas operativos de tiempo real no son sistemas operativos completos. Estan compuestos, basicamente, por un gestor de tareas apropiativo y de mecanismos para la sincronizacion y comunicacion entre las mismas. Este reducido conjunto de caracteristicas esta directamente relacionado con el hardware donde se ejecuta: los RTOS son utilizados mayormente en plataformas de computo basado en microcontroladores de bajos recursos para automatizacion y control.

En los ultimos anos, con el crecimiento exponencial de desarrollo de dispositivos IOT, su uso ha crecido en ambientes de sensores conectados a Internet, tambien controlados generalmente por microcontroladores. Algunos ejemplos de aplicaciones embebidas típicas desarrolladas utilizando un RTOS son: 
\begin{itemize}
	\item Robots moviles, : control de motores, lectura de encoders opticos, calculo de odometria.
	\item Sistemas de adquisicion de datos: muestreo periodico de sensores con restricciones temporales.
	\item Controladores de procesos industriales: automatizacion con respuesta garantizada ante eventos.
	\item Dispositivos IOT: sensores conectados que reportan valores dentro de tiempos limite.
	\item Sistemas de navegacion: fusion de datos de magnetometro, acelerometro y encoders con plazos estrictos.
\end{itemize}

\section{El limite del bucle infinito}

Para apreciar por qué es muy comun utilizar un RTOS al desarrollar un sistema embebido
comenzaremos con un ejemplo típico.
		Antes de introducir un RTOS, muchos sistemas embebidos de baja complejidad se implementan con un esquema simple: un bucle infinito principal (superloop) que atiende la logica de la aplicacion, combinado con rutinas de servicio de interrupcion (ISR) para manejar eventos asincronicos. Este esquema es facil de entender y funciona bien cuando existe un unico requisito de tiempo de respuesta o las tareas tienen plazos muy similares. Pero en cuanto aparecen dos o mas requisitos temporales diferentes, el diseno se vuelve rapidamente inviable.

La razon de fondo es sencilla: en un bucle infinito, todas las tareas comparten el mismo hilo de ejecucion. No existe ningun mecanismo que garantice que una tarea critica sea atendida a tiempo si otra tarea menos critica esta en ese momento ocupando la CPU. La ISR puede despertar un flag o actualizar una variable, pero si el bucle principal esta en medio de una operacion larga, el procesamiento del evento llegara tarde.

\subsection{Ejemplo: robot movil con odometria y temperatura}

Suponga que se debe desarrollar un sistema embebido robotico que navega y controla la temperatura en diferentes ambientes.
El robot realiza dos funciones principales:

\begin{itemize}
	\item Temperatura: el sistema debe leer varios sensores de temperatura. Si al procesar esos valores se exede un umbral se activa una alarma. Esta función debe realizarse al menos \textbf{una vez cada 10ms}. Al implementar esta funcionalidad se evalúa y se demuestra que esta funcionalidad se logra en \textbf{3ms} utilizando este microcontrolador.
	\item Odometria: el sistema debe leer encoders opticos que permiten calcular la velocidad de cada rueda del robot, y con esta información, calcular la posicion del robot \textbf{una vez cada 5ms} al menos. Si el calculo se realiza luego de 5ms el error de posicionamiento supera el margen aceptable. Al implementar la funcionalidad se evalúa y se demuestra que el computo demora \textbf{3ms} en este microcontrolador. \textbf{Esta tarea se define como la de mayor importancia}.
\end{itemize}

\subsubsection{Intento con bucle infinito e interrupciones}

Una primera implementación tradicional es utilizar dos timers por hardware para activar las diferentes tareas a tiempo, mediante interrupciones (cada 10ms y cada 5ms). Luego, procesar todo en el bucle principal de main.

\begin{verbatim}
// ---- Intento con bucle + ISR: PROBLEMATICO ----

ISR(TIMER1_OVF_vect) {        // interrupcion cada 10 ms
  temperaturas = leer_sensores_temp();
  tarea_temp = 1;
}

ISR(TIMER0_OVF_vect) {        // interrupcion cada 5 ms
  velocidad = leer_encoder();
  tarea_odom = 1;
}

int main(void) {
  for(;;) {
      if (tarea_temp) {	  /* temperatura */
          tarea_temp = 0;
          PELIGRO = procesar(temperaturas);    // demora 3ms
                                               // PROBLEMA !!
          if (PELIGRO)
              alarma_ON();	
      }

      if (tarea_odom) {	  /* odometria */
          tarea_odom = 0; 		
          calcular_odometria(velocidad);      // demora 3ms
      }
  }
}
\end{verbatim}

\subsubsection*{Donde falla este diseño}

\begin{itemize}
	\item procesar(temperaturas) demora 3ms, y debe finalizar este cálculo en menos de 10ms desde que se activó la tarea. 

	\item calcular\_odometria() demora 3ms, y debe finalizar este cálculo en menos de 5ms desde que se activó la tarea. 
\end{itemize}

El problema con la solución tradicional es que si las tareas se activan en algun momento casi simultaneamente,  con tarea\_temp en primer lugar, entonces 
tarea\_odom finalizará a los 6ms, lo que producirá como resultado que la \textbf{tarea crítica superó el tiempo límite}.
Tal incumplimento del plazo producirá un error irrecuperable. En este ejemplo, el robot no logra calcular su posición a tiempo, y probablemente comenzará a malfuncionar (chocando o realizando movimientos fuera de lugar).


% "La tarea no cumplió con el plazo."
% "El sistema superó el tiempo límite."
% "El proceso no alcanzó a terminar dentro del tiempo establecido."
% "Se venció el tiempo antes de que terminara el cómputo."
% "La tarea excedió el tiempo máximo permitido."
% "No se logró completar el cálculo dentro de la ventana de tiempo."

Una alternativa que inmediatamente muchos estudiantes intentan proponer es modificar las rutinas de atencón de interrupciones, y colocar dentro el procesamiento de cada tarea. De esta manera, si se está procesando las temperaturas, y aparece una interrupción de odometría, se ejecutará inmediatamente la rutina de atención de interrupciones calculando la odometría a tiempo.
Esta solución, aparentemente apropiada en principio, tiene el inconveniente que las rutinas de atención de interrupciones toman mucho tiempo en ejecutarse. 
Como se vió en módulos teóricos previos, una rutina de atención de interrupciones debe activarse, realizar alguna acción lo más rápido posible y finalizar. No es aceptable tener a estas rutinas en ejecución por varios ms. La principal razón es que otras interrupciones no serán atendidas (mientras se ejecuta una rutina de este tipo las interrupciones están desactivadas). En caso de activar las interrupciones dentro de rutinas críticas, para no perder otras interrupciones, aún sigue siendo un gran problema mantener este código. Casi imposible ampliarlo con mayores funcionalidades por rutinas de atención de interrupciones anidadas.



\subsubsection{Solucion con RTOS y prioridades}

Utilizando como herramenta de apoyo un RTOS, cada requisito se convierte en una tarea independiente con su propia prioridad. El planificador siempre asigna la CPU a la tarea de mayor prioridad lista para ejecutarse. 
Esto es, independientemente de lo que este haciendo alguna otra tarea de menor prioridad, 
cuando utiliza la CPU, si una tarea de mayor prioridad se activa, el planificador de CPU del RTOS realiza un cambio de contexto y pone a ejecutar la tarea de mayor prioridad.

\begin{verbatim}
// ---- Solucion utilizando un RTOS ----

// PRIORIDAD 15 — 10 ms
process tarea_temperatura(void) {
   for (;;) {
        temperaturas = leer_sensores_temp();
        PELIGRO = procesar(temperaturas);     // demora 3ms
                        
        if (PELIGRO)
              alarma_ON();	

        sleepms(10);
    }
}

// PRIORIDAD 40 — tarea critica: finalizar antes de 5 ms
process tarea_odometria(void) {
   for (;;) {
        velocidad = leer_encoder();
        calcular_odometria(velocidad);      // demora 3ms
	
        sleepms(5);   
    }         
}

void main(void) {
    resume(create(tarea_temperatura,256, 15, "temp", 0));
    resume(create(tarea_odometria,  256, 40, "odom", 0));
}
\end{verbatim}

Aqui, si tarea\_temperatura se está ejecutando, y en ese
instante vence el periodo de espera de 5ms de la tarea\_odometria, 
el planificador interrumpe a 
tarea\_temperatura y le da la CPU a tarea\_odometria. 
Cuando esta termina, el RTOS bloquea a tarea\_odometria por otros 5ms. 

Es decir, sleepms() no hace una espera activa. sleepms()
es una llamada al sistema del RTOS
que pone a la tarea en estado SLEEPING. En ese estado, la tarea nunca
consume CPU. Cuando sucedan exactamente 5ms, el RTOS pone a la tarea\_odometria en estado de LISTO. En este punto, inmediatamente la coloca en estado EJECUTANDO y realiza un cambio de contexto para que utilice la CPU, 
ya que al ingresar al estado de LISTO, y ser de la mas alta prioridad, 
el RTOS asegura su ejecución inmediata.
El plazo de 5 ms se cumple siempre, sin necesidad de calcular 
manualmente cuanto tarda cada funcion.

Un lector atento observará que poner sleepms(5); no sería tan adecuado.
Si la tarea demora 3ms en ejecutarse, entonces la tarea se estaría
ejecutando cada 8ms (y no cada 5ms). sleepms(2) sería mas apropiado.
Otra solución mas eficiente podría ser utilizando un timer y un semáforo
provisto por el RTOS para activar la tarea, o tener una tarea que haga
de temporizador con igual mecanismo.

El lector puede además, utilizando este ejemplo, observar lo sencillo
que sería agregar mayores funcionalidades. Si el sistema requiere
otras tareas, simplemente se las desarrolla de manera independiente, como 
en este caso. Utilizando las prioridades adecuadas se puede modelar
el comportamiento de manera predecible, y cumplir con los tiempos
de respuesta necesarios para cada funcionalidad.

\subsection{Ejemplo de un sistema con eventos asincronicos}
Ahora consideremos un sistema que ademas debe atender comandos que llegan por UART desde una computadora. Los comandos no llegan en forma periodica: pueden llegar en cualquier momento, y el sistema debe responder dentro de 20 ms desde la recepcion. Si no responde a tiempo, la comunicacion se pierde.
Agregamos este requisito al escenario anterior:
    • Requisito A — Odometria: cada 5 ms. Prioridad maxima.
    • Requisito B — Comandos UART: responder en menos de 20 ms desde la llegada del byte. Prioridad alta.
    • Requisito C — Temperatura: cada 100 ms. Prioridad baja.

Intento con bucle infinito + ISR

\begin{verbatim}
// ---- Tres requisitos con bucle + ISR: AUN MAS PROBLEMATICO ----

volatile bool flag_odom  = false;
volatile bool flag_uart  = false;
volatile bool flag_temp  = false;
volatile char cmd_byte   = 0;

ISR(TIMER0_OVF_vect)  { flag_odom = true; }           // cada 5 ms
ISR(USART_RX_vect)    { cmd_byte = UDR0;
                         flag_uart = true; }           // asincrono
ISR(TIMER1_OVF_vect)  { flag_temp = true; }           // cada 100 ms

int main(void) {
    sei();
    while (1) {
        // El orden de atencion lo define el programador, no las prioridades reales.
        // Si odometria tarda 6 ms y justo llega un comando UART, perdemos el plazo.
        if (flag_odom) { flag_odom = false; calcular_odometria(); }
        if (flag_uart) { flag_uart = false; procesar_comando(cmd_byte); }
        if (flag_temp) { flag_temp = false; verificar_temperatura(); }
    }
}
\end{verbatim}

El problema de fondo
El orden de atencion de los flags en el bucle lo decide el programador en tiempo de escritura del codigo, no las prioridades reales del sistema en tiempo de ejecucion. Si calcular\_odometria() tarda mas de lo esperado y durante esa ejecucion llega un byte UART, el flag\_uart se activa pero no sera atendido hasta que el bucle vuelva a chequearlo. Si para entonces han pasado mas de 20 ms, el plazo se perdio.
Este problema no tiene solucion limpia dentro del esquema de bucle + ISR. Se pueden hacer particiones de tiempo manuales, dividir las funciones en "slices" mas pequenos, o usar maquinas de estado que se ejecutan por partes... pero toda esa complejidad es exactamente lo que un RTOS resuelve de forma sistematica y formal.

Solucion con RTOS: eventos que despiertan tareas
Con un RTOS, la ISR de UART simplemente despierta a la tarea de comandos via un semaforo. La tarea no consume CPU mientras espera; en cuanto llega el byte, el planificador la activa y, si su prioridad es mayor que la tarea en ejecucion, la pone a correr inmediatamente:

\begin{verbatim}
// ---- Solucion con XINU (AVR) ----

sid32 sem_uart;   // semaforo global, inicializado en 0
volatile char cmd_byte = 0;

// ISR: solo guarda el dato y despierta a la tarea consumidora
ISR(USART_RX_vect) {
    cmd_byte = UDR0;
    signal(sem_uart);   // despierta tarea_comandos si esta bloqueada
}

// PRIORIDAD 40 — critica: 5 ms
process tarea_odometria(void) {
    while (TRUE) {
        calcular_odometria(leer_encoder());
        sleep(5);
    }
}

// PRIORIDAD 30 — alta: responde en cuanto llega el byte (< 20 ms garantizado)
process tarea_comandos(void) {
    while (TRUE) {
        wait(sem_uart);          // bloqueada hasta que ISR haga signal()
        procesar_comando(cmd_byte);
    }
}

// PRIORIDAD 10 — baja: 100 ms
process tarea_temperatura(void) {
    while (TRUE) {
        if (leer_sensor_temp() > UMBRAL) encender_LED();
        else                             apagar_LED();
        sleep(100);
    }
}

void main(void) {
    sem_uart = screate(0);
    resume(create(tarea_odometria,  256, 40, "odom", 0));
    resume(create(tarea_comandos,   256, 30, "cmd",  0));
    resume(create(tarea_temperatura,256, 10, "temp", 0));
}
\end{verbatim}

Cuando la ISR ejecuta signal(sem\_uart), el planificador de XINU despierta a tarea\_comandos. Como su prioridad (30) es mayor que la de tarea\_temperatura (10), si esta ultima estaba ejecutando, es interrumpida inmediatamente. El tiempo entre la llegada del byte y el inicio de procesar\_comando() es del orden de microsegundos: el plazo de 20 ms se cumple con amplio margen.

Escenario	Con bucle + ISR	Con RTOS y prioridades
1 tarea periodica	Funciona bien	Funciona bien (mas simple de escribir)
2 tareas con plazos distintos	Requiere analisis manual de tiempos	El planificador garantiza los plazos
3+ tareas con plazos distintos	Muy dificil, propenso a errores	Cada tarea tiene su prioridad; sistematico
Evento asincrono (UART, sensor)	El flag puede llegar tarde	ISR + signal(): respuesta inmediata
Agregar una nueva tarea	Redisenar el bucle y recalcular tiempos	Crear la tarea con la prioridad adecuada

3. Patrones tipicos de uso de tareas en un RTOS
Mas alla de los ejemplos anteriores, en la practica, casi todas las tareas de un sistema embebido con RTOS caen en uno de los siguientes tres patrones. Conocerlos permite disenar el sistema de forma ordenada desde el principio.

\subsection{ Tarea periodica (hace algo y duerme)}
La tarea realiza su trabajo y luego duerme por un periodo fijo. Al despertar, vuelve a ejecutarse. Este patron es el mas directo para requisitos como muestreo de sensores, control PID, actualizacion de displays, o calculo de odometria. El RTOS garantiza que, si la tarea tiene la prioridad adecuada, sera activada a tiempo independientemente de lo que hagan las demas.

\begin{verbatim}
process tarea_control_pid(void) {
    while (TRUE) {
        int error = setpoint - leer_sensor_velocidad();
        int salida = calcular_pid(error);
        aplicar_pwm(salida);
        sleep(10);   // periodo de 10 ms
    }
}
\end{verbatim}

\subsection{ Tarea que espera un evento (semaforo o senal)}
La tarea permanece bloqueada hasta que ocurre un evento externo. Mientras espera, no consume CPU: el planificador la ignora y ejecuta otras tareas. En cuanto el evento ocurre y alguien hace signal() sobre su semaforo, la tarea se desbloquea y, si su prioridad es la mas alta entre las listas, el planificador le da la CPU inmediatamente.
Este patron es ideal cuando la tarea no debe estar activa continuamente, sino solo cuando hay trabajo real que hacer: llegada de un comando por UART, deteccion de un evento de hardware, finalizacion de una conversion ADC, etc.

\begin{verbatim}
sid32 sem_conversion_adc;

ISR(ADC_vect) {
    resultado_adc = ADCW;
    signal(sem_conversion_adc);   // conversion terminada: despierta la tarea
}

process tarea_procesar_adc(void) {
    while (TRUE) {
        wait(sem_conversion_adc);   // bloqueada hasta que ISR haga signal()
        procesar_muestra(resultado_adc);
        iniciar_nueva_conversion();
    }
}
\end{verbatim}

\subsection{ ISR como productora, tarea como consumidora}
Los datos llegan por interrupcion en forma rapida y frecuente. La ISR no puede procesar: solo debe depositar el dato en un buffer y despertar a la tarea consumidora. La tarea consumidora procesa los datos a su ritmo, sin riesgo de perder interrupciones porque el buffer actua como cola.
Este patron es especialmente util para UART, SPI, I2C, o cualquier periferico que genera datos a alta velocidad. La ISR es minima y rapida; el procesamiento pesado queda en la tarea.

\begin{verbatim}
#define BUF_SIZE 16

volatile char buffer[BUF_SIZE];
volatile int  buf_in    = 0;
volatile int  buf_count = 0;
sid32         sem_datos;   // inicializado en 0

ISR(USART_RX_vect) {
    if (buf_count < BUF_SIZE) {
        buffer[buf_in] = UDR0;
        buf_in = (buf_in + 1) % BUF_SIZE;
        buf_count++;
    }
    signal(sem_datos);   // despierta al consumidor
}

process tarea_consumidor_uart(void) {
    int out = 0;
    while (TRUE) {
        wait(sem_datos);             // espera hasta que haya dato
        char d = buffer[out];
        out = (out + 1) % BUF_SIZE;
        buf_count--;
        procesar_caracter(d);
    }
}
\end{verbatim}

Resumen de patrones
Periodica:  sleep(N ms) al final del ciclo. Ideal para control, muestreo, actualizacion periodica.
Evento:  wait(sem) hasta que alguien haga signal(). Ideal para responder a eventos externos sin consumir CPU.
ISR + consumidor:  ISR deposita en buffer y hace signal(); tarea procesa a su ritmo. Ideal para perifericos de alta velocidad (UART, ADC, SPI).
Estos tres patrones, combinados con prioridades bien asignadas, cubren la gran mayoria de los sistemas embebidos reales.

\section{XINU en Arduino (AVR atmega328p)}
XINU es un sistema operativo academico desarrollado originalmente por Douglas Comer en la Universidad de Purdue, a fines de los anos 70. Desde entonces, ha sido re-escrito y portado a una gran variedad de plataformas de hardware y arquitecturas, y es utilizado, principalmente, como herramienta de investigacion y educacion.
A pesar de que XINU comparte algunos conceptos y alguna terminologia de componentes con UNIX, el diseno interno de XINU difiere completamente de UNIX. XINU es un sistema operativo con un diseno elegante, cuya estructura interna esta conformada por una jerarquia multinivel de componentes esenciales, pero con las cuales puede luego desarrollarse otras mas complejas. Tiene soporte para la creacion dinamica de procesos, reserva dinamica de memoria, sistemas de archivos, soporte de red, y funciones de E/S independiente de los dispositivos. Tambien ofrece servicios de comunicacion y sincronizacion entre procesos.
XINU no implementa memoria virtual, ni tiene soporte para hardware de paginacion. En tiempo de ejecucion, la imagen de XINU con un conjunto de aplicaciones se carga completamente en memoria RAM, utilizando un unico espacio de direcciones para todos los procesos. Esto significa que los procesos comparten la memoria, aunque XINU gestiona para cada proceso su propia pila.
El codigo fuente del kernel de XINU contiene aproximadamente 2500 lineas de codigo en lenguaje C, lo que lo hace completamente explorable en el contexto de una materia de grado. Esta es una de sus principales virtudes academicas: no solo se puede usar como RTOS, sino tambien entender, modificar y extender su funcionamiento interno.

XINU para AVR (Arduino Nano)
El port de XINU para arquitectura AVR corre sobre el microcontrolador atmega328p: arquitectura Harvard, 32 KB de memoria flash y 2 KB de RAM. Fue evaluado exitosamente en:
  - La materia Programacion de Sistemas Embebidos (Lic. en Ciencias de la Computacion, UNComa).
  - Un trabajo experimental de tesis de grado que requirio un RTOS para controlar un robot movil con magnetometro y encoders opticos, cuyos eventos asincronicos requerian reportar valores en tiempos limite para que el calculo de odometria tuviese un error en los margenes esperados.
Codigo fuente disponible en:  http://se.fi.uncoma.edu.ar/xinu-avr/

\subsection{ Creacion y finalizacion de procesos}
En XINU, un proceso es una funcion C que se ejecuta de forma concurrente con las demas. La funcion create() crea el proceso en estado suspendido; resume() lo coloca en estado "listo para ejecutar". La combinacion de ambas es la forma estandar de lanzar un proceso. create() retorna un PID (Process ID) que permite referirse al proceso en llamadas posteriores.
La firma de create() es:
\begin{verbatim}
// create(funcion, tam_pila_bytes, prioridad, nombre, nro_args, args...)
pid32 create(void *funcion, uint16 stksize, pri8 priority,
             char *name, uint8 nargs, ...);
Un ejemplo completo de creacion, obtencion de PID y finalizacion:
#include <xinu.h>

process mi_tarea(int param);   // prototipo

void main(void) {
    pid32 pid;

    pid = create(mi_tarea, 256, 20, "mi_tarea", 1, 42);

    if (pid == SYSERR) {
        kprintf("Error al crear el proceso\n");
        return;
    }

    kprintf("Proceso creado con PID = %d\n", pid);
    resume(pid);   // lo coloca en estado LISTO para ejecutar
}

process mi_tarea(int param) {
    kprintf("Hola desde mi_tarea, param = %d\n", param);
    // Al retornar de la funcion, el proceso finaliza automaticamente.
    // Tambien se puede terminar desde afuera con kill(pid).
}

Otras primitivas de gestion de ciclo de vida:
pid32 pid_actual = getpid();       // PID del proceso que ejecuta esta linea

kill(pid);                          // termina un proceso por su PID

suspend(pid);                       // suspende un proceso (queda en estado SUSPENDED)
resume(pid);                        // reanuda un proceso suspendido

chprio(pid, nueva_prioridad);       // cambia la prioridad de un proceso en ejecucion
\end{verbatim}

4.2. Procesos periodicos con sleep()
La funcion sleep(ms) bloquea el proceso durante la cantidad de milisegundos indicada. Durante ese tiempo, el planificador cede la CPU a otros procesos listos. Cuando el tiempo vence, el proceso pasa al estado LISTO y sera ejecutado cuando el planificador lo decida segun su prioridad.
Combinada con un bucle infinito dentro de la tarea, sleep() implementa el patron de tarea periodica de forma directa. A diferencia del bucle + ISR, aqui cada tarea tiene su propio plazo, y el planificador garantiza que la de mayor prioridad siempre se activa primero:
\begin{verbatim}
// Dos tareas periodicas con distintas restricciones temporales

// PRIORIDAD 40 — critica: error de odometria inaceptable si se retrasa
process tarea_odometria(void) {
    while (TRUE) {
        calcular_odometria(leer_encoder());
        sleep(5);    // 5 ms: restriccion critica
    }
}

// PRIORIDAD 15 — menos critica: tolerancia de decenas de ms
process tarea_temperatura(void) {
    while (TRUE) {
        int t = leer_sensor_temp();
        if (t > UMBRAL) encender_LED();
        else             apagar_LED();
        sleep(100);  // 100 ms
    }
}

void main(void) {
    resume(create(tarea_odometria,  256, 40, "odom", 0));
    resume(create(tarea_temperatura,256, 15, "temp", 0));
}
\end{verbatim}

Por que esto funciona y el bucle no
Si tarea\_temperatura esta ejecutando y vence el periodo de tarea\_odometria, el planificador interrumpe inmediatamente a tarea\_temperatura y le da la CPU a tarea\_odometria. tarea\_temperatura retoma exactamente donde fue interrumpida cuando tarea\_odometria termina su ciclo. El programador no necesita calcular cuanto tarda cada funcion: el planificador lo gestiona de forma sistematica en tiempo de ejecucion.

\subsection{ Sincronizacion con semaforos: productor y consumidor}
Los semaforos son el mecanismo principal de sincronizacion en XINU. Conceptualmente, un semaforo es un contador entero con dos operaciones atomicas:
    • wait(sem): si el contador es mayor que cero, lo decrementa y continua. Si es cero, bloquea el proceso hasta que otro haga signal().
    • signal(sem): incrementa el contador. Si hay procesos bloqueados esperando, desbloquea al primero (segun prioridad).
La funcion screate(n) crea un semaforo con valor inicial n. El siguiente ejemplo muestra el patron productor/consumidor completo con buffer circular:

\begin{verbatim}
#include <xinu.h>

#define BUF_SIZE 8

sid32 sem_lleno;          // cuantos items hay en el buffer (inicia en 0)
sid32 sem_vacio;          // cuantos huecos libres hay (inicia en BUF_SIZE)
int   buffer[BUF_SIZE];
int   in = 0, out = 0;

process productor(void) {
    int i = 0;
    while (TRUE) {
        wait(sem_vacio);              // si buffer lleno, se bloquea
        buffer[in] = i++;
        in = (in + 1) % BUF_SIZE;
        signal(sem_lleno);            // avisa que hay un item nuevo
        sleep(10);                    // produce cada 10 ms
    }
}

process consumidor(void) {
    while (TRUE) {
        wait(sem_lleno);              // si buffer vacio, se bloquea
        int dato = buffer[out];
        out = (out + 1) % BUF_SIZE;
        signal(sem_vacio);            // libera un hueco
        kprintf("Consumido: %d\n", dato);
    }
}

void main(void) {
    sem_lleno = screate(0);           // buffer inicialmente vacio
    sem_vacio = screate(BUF_SIZE);    // todos los huecos libres

    resume(create(productor,  256, 20, "prod", 0));
    resume(create(consumidor, 256, 20, "cons", 0));
}
\end{verbatim}

Que garantizan los semaforos aqui?
El consumidor nunca lee del buffer vacio: sem\_lleno comienza en 0 y lo bloquea hasta que el productor haga signal().
El productor nunca escribe en un buffer lleno: sem\_vacio lo bloquea cuando llega a 0.
No se necesita deshabilitar interrupciones ni usar flags volatiles con acceso no protegido: el RTOS maneja la sincronizacion de forma segura y predecible. Cualquier violacion de estas reglas quedaria capturada por el semaforo antes de ocurrir.

\subsection{ Ejemplo completo: tres tareas con prioridades distintas}
El siguiente ejemplo integra los tres patrones en un sistema realista: tarea periodica critica, tarea despertada por evento asincrono, y tarea periodica de baja prioridad. Corresponde al tipo de diseno utilizado en el trabajo de tesis de grado mencionado anteriormente:

\begin{verbatim}
#include <xinu.h>

sid32 sem_uart;
volatile char cmd_byte;

// ISR: rapida, solo guarda dato y despierta la tarea
ISR(USART_RX_vect) {
    cmd_byte = UDR0;
    signal(sem_uart);
}

// PRIORIDAD 40 — critica: odometria, 5 ms
process tarea_encoder(void) {
    while (TRUE) {
        calcular_odometria(leer_encoder());
        sleep(5);
    }
}

// PRIORIDAD 30 — alta: responde a comandos UART en < 20 ms
process tarea_comandos(void) {
    while (TRUE) {
        wait(sem_uart);
        procesar_comando(cmd_byte);
    }
}

// PRIORIDAD 10 — baja: display LCD, 500 ms
process tarea_display(void) {
    while (TRUE) {
        actualizar_lcd();
        sleep(500);
    }
}

void main(void) {
    sem_uart = screate(0);
    resume(create(tarea_encoder,  256, 40, "encoder",  0));
    resume(create(tarea_comandos, 256, 30, "comandos", 0));
    resume(create(tarea_display,  256, 10, "display",  0));
}
\end{verbatim}

El planificador garantiza el siguiente comportamiento en tiempo de ejecucion: tarea\_encoder siempre se activa primero cuando su periodo de 5 ms vence, sin importar que esten haciendo las otras. tarea\_comandos se activa en microsegundos luego de que llega un byte UART, con tiempo de sobra para cumplir el plazo de 20 ms. tarea\_display solo ocupa la CPU cuando las otras dos estan dormidas o bloqueadas. Todo esto sin que el programador tenga que calcular manualmente cuanto tarda cada funcion ni en que orden se ejecutan.

5. Sintesis y proximos pasos
A lo largo de este tutorial hemos visto que un sistema de tiempo real necesita garantizar no solo la correccion de sus resultados, sino tambien el cumplimiento de plazos. El uso de un RTOS es una de las formas mas comunes de lograr esto, aunque no la unica. Su ventaja central es que permite disenar el sistema como un conjunto de tareas independientes, cada una con su propia prioridad, y delegar al planificador la responsabilidad de garantizar que la tarea mas urgente siempre ocupe la CPU.
El esquema de bucle infinito + ISR se vuelve inviable a medida que crecen la cantidad de tareas y la diversidad de sus plazos. No existe una forma sistematica de garantizar todos los plazos dentro de ese esquema. Un RTOS resuelve este problema de forma formal y escalable: agregar una nueva tarea es tan simple como llamar a create() con la prioridad adecuada.
XINU para AVR (Arduino) es una herramienta ideal para aprender estos conceptos: su codigo fuente es pequeno y explorable, su API es sencilla y uniforme, y corre en el hardware accesible de una placa Arduino Nano. Los ejemplos de este tutorial pueden ejecutarse directamente en ese hardware.

Para seguir explorando
    1. Codigo fuente de XINU-AVR: http://se.fi.uncoma.edu.ar/xinu-avr/
    2. Libro de referencia: D. Comer, "Operating System Design — The Xinu Approach", 2da ed., CRC Press, 2015.
    3. Ejercicio 1: ampliar el ejemplo del productor/consumidor para que la ISR de UART sea el productor y una tarea XINU sea el consumidor que procesa comandos. Verificar con un terminal serie.
    4. Ejercicio 2: disenar un sistema con 3 tareas de distintas prioridades y medir con un osciloscopio (o con GPIO toggling) los tiempos de activacion. Verificar que la tarea de mayor prioridad siempre cumple su plazo.
    5. Ejercicio 3: implementar el patron ISR + consumidor para la recepcion por I2C de un sensor IMU, y comparar el jitter en el tiempo de procesamiento contra un diseno con bucle + ISR.



























La interfaz entre un procesador embebido y un dispositivo periférico
es un conjunto de registros de estado y control. Estos registros son parte
del hardware del periférico, y sus ubicaciones, sus tamaños (en bits) y el
significado individual de cada bit en cada registro son característicos 
de cada dispositivo. Por ejemplo, 
el significado de los bits de los registros de un dispositivo serial son 
muy diferentes a los
de los contadores o relojes. En esta sección se encuentra explicado la 
manera de manipular el contenido de los registros de estado y control
directamente desde programas en C/C++.

Dependiendo del diseño del procesador y la placa, los dispositivos periféricos
se encuentran localizados en el espacio de memoria del procesador o dentro
del espacio de E/S. De hecho, es común en sistemas embebidos que se incluyan
periféricos de ambos tipos. Se los denomina \textit{periféricos mapeados en memoria}
y \textit{periféricos mapeados en E/S} respectivamente (estos últimos tambien son 
denominados como \textit{aislados}). De los dos tipos, los periféricos
mapeados en memoria son los más fáciles de utilizar y por lo tanto su popularidad
aumenta.

Los registros de estado y control de periféricos mapeados en memoria se pueden
parecer a variables comunes en un programa. Por ejemplo, es posible 
declarar un puntero en C a un registro o bloque de registros, 
y establecer
la dirección explícitamente (valor del puntero). Supongamos el registro
\texttt{puerto\_b} utilizado en el capítulo 2, el cual es mapeado en memoria en la dirección 
física 0x25. La función \texttt{led\_toggle} estudiada en ese capitulo puede ser 
escrita completamente en C, como se muestra a continuación.

\begin{verbatim}
volatile unsigned char * puerto_b = (unsigned char *) 0x25;

void led_toggle(void)
{
    *puerto_b ^= LED_ROJO;    /* Read, xor, and modify. */
}    /* led_toggle() */
\end{verbatim}

Aquí se declara un puntero a un registro de 8 bits y explícitamente se lo
inicializa con  la dirección 0x25. A partir de este momento el puntero al 
registro
trabaja como cualquier puntero a una variable de tipo char de 8 bits.


\section{Referencias}

Michael Barr. Programming Embedded Systems in C and C++ 1st Edition. ISBN-13: 978-1565923546
ISBN-10: 1565923545. O'Reilly Media; 1 edition (February 9, 1999)


\section {Licencia y notas de la traducción}

Rafael Ignacio Zurita ({\texttt{rafa@fi.uncoma.edu.ar}) 2016-2020

Este apunte es una traducción del libro de referencia, para
ser utilizado como apunte en la materia de grado
'Programación de Sistemas Embebidos' de la Facultad de Informática,
Universidad Nacional del Comahue.
Se han realizado modificaciones
al contenido para aclarar ciertos detalles o agregar secciones nuevas.
 Tambien se han
modificado todos los archivos fuentes de código, ya que fueron
portados a la plataforma Atmel AVR atmega328.


\begin{center}
\begin{tabularx}{\textwidth}{|X|}
%\hline
\rowcolor{aliceblue}
\textit{Esta es una obra derivada del libro de referencia que ha obtenido permiso de O'Reilly para su publicación en PEDCO.}\\
%\hline
\end{tabularx}
\end{center}



\subsubsection {Permiso de publicación}

\begin{footnotesize}
\begin{verbatim}
From: Teri Finn <teri@oreilly.com>
Date: Mon, 11 Apr 2016 15:48:58 -0700
Message-ID: <CAOgscZ92-G9iT+=+nnuft4LCoHoe5o8SYPfVEKnS6AKXh8TLyQ@mail.gmail.com>
Subject: Re: Question about permission to translate some parts
To: Rafael Ignacio Zurita <rafa@fi.uncoma.edu.ar>
Content-Type: multipart/alternative; boundary=94eb2c094eeedeba0005303d5a31

--94eb2c094eeedeba0005303d5a31
Content-Type: text/plain; charset=UTF-8

Hi Rafael,

Thank you for providing this further information.  O'Reilly Media is happy
to grant you permission to translation the content you have proposed for
posting to the Moddle site. We're happy you find this information helpful
to the students.

Best regards,

Teri Finn
O'Reilly Media, Inc.

On Thu, Apr 7, 2016 at 10:56 AM, Rafael Ignacio Zurita <
rafa@fi.uncoma.edu.ar> wrote:

[snip]
\end{verbatim}
\end{footnotesize}



%\end{otherlanguage}
% } % foreign language


\end{document}








